\documentclass[12pt]{article}

\usepackage[a4paper, total={6in, 8in}]{geometry}
\usepackage{textcomp}

\begin{document}
\setlength{\parindent}{0pt}

\def \t {\textrightarrow}
\def \v {\vspace{1em}}
\def \bi {\begin{itemize}}
\def \ei {\end{itemize}}
\def \s[#1] {\section*{#1}}
\def \ss[#1] {\subsection*{#1}}
\def \sss[#1] {\subsubsection*{#1}}

\s[Ungaretti]
\ss[San Martino del Carso]
Paesaggistica straziata dalla macchina bellica \t il dolore nelle mani, qua si incarna in un paesaggio distrutto \t il neorealismo parte da questa immagine (es. Pavese) \\
Muro = emblema limitante \t non è il muro, ma è un brandello \\ 
Presenti molte negazioni nel testo, e suo dei partitivi \\
Verso 9 \t ma cambia il registro in modo avversativo: questo crea spezzatura nel verso, anche se il verso è continuativo (punteggiatura praticamente inesistente) \\
Il "ma" discrima l'esterno dall'interno, la desolazione del paesaggio da quella dell'animo \\
La lente è quella del cuore \t Palazzeschi mette una lente nel cuore \t qua il dolore è nel suo cuore \\
Straziato \t orologio a rote, si dice che lacera il giorno l'orologio \t stessa cosa qua: il cuore è lacerato, straziato dal dolore \t il paesaggio e il cuore sono in una mimesis totale: devastazione interna ed esterna \\
Connotazioni stilistiche: nuova armonia parola scritta, lettura attenta suggerisce la volonta dell autore di arrivare con il suo dolore al lettore (l'immagine che si legge viene resa palpitante attraverso elementi cronologici e geografici) \\
Anno e luogo sono presenti, e questo conferisce a ogni testo come un tassello \t ogni composizione è come se fosse la pagine di un diario \\
In questa lirica abbiamo quindi la congestione che penetra nelle mani \t la nelle mani, qua la lettura è tutta sua \t è il dolore universale, che si fa personale \\

\ss[Sono una creatura]
È il testo dell'anafora \t ripetuto "cosi", e "come" è una analogia esplicita \\
Realta che mostra il dolore, che dal paesaggio si trasponde nel suo cuore e si traduce in pianto \\
I versi sono costituiti da una sola espressione \t un climax ascendente, che va dal verso 3 al 8 \\
1916, anche questo in "Vita di un uomo" \\
Versi con una sola parola, non presenza di punteggiatura \t assimetria tra prima e ultime strofe, che sono molto + corte \\
La morte si sconta vivendo \t espresso qua il senso della vita \t il senso religioso mette davanti alla vita, che è da vedersi come preparazione di altro \t la morte è vissuta in ogni attimo di dolore, che è penetrato nel suo animo e si traduce nel suo pianto \\
San Michele è il monte, e pietra è metonimia (parte della montagna) \t aridità, assenza di vegetazione, che ricorda i pendii della Ginestra \\
Prosciugata e disanimata possono essere visti come particip erfetti \t tempo che chiude totalmente il cerchio, è un azione completamente conclusa \\
Qua illustra come la morte è presente in questa vita, che non è neanche + vita \t ha reso coloro che hanno sperimentato la morte, tali da essere queli inanimati da essere pietre \\
Questo è il dolore non urlato del grido di Munch \t è il dolore muto \t questo è il dolore vissuto metafisicamente \\
La morte, come unica prospettiva 

\ss[Vanità]
1917 \\
Vanità \t nell'inferno sono gli ignavi \t Petrarca utilizzava la parola pura, poi si sporca con d'annunzio \t ma approda qui, dove vanita assume significato esistenziale \\
La poesia degli anni fino al 30 è illuminazione, fulminea \t anche qua versi monolitici, rinnovato equilibrio, e nella costruzione esistenziale che non è dei futuristi (ma di un autore che perocrre anche il futurismo) \\
E anche una profonda espressivita, data negli ultimi vesi, che apre un varco e una finestra nell'ambito montaliano \t l'immagine riflessa è molto presente, come "Cigola la carrucola nel pozzo" \\
La guerra non ha solo devastato, ma ha anche distrutto la consapevolezza di se \t l'uomo è stato capace di tanta violenza \t non c'e immagine della coscienza di zeno, secono cui distruggere è input per ricostruire umanita nuova \\
L'immenso e la piccolezza dell'uomo \t le macerie cozzano con il limpido stupore dell'immensita \t il limite è nelle macerie \\
Anche l'acqua è raggiunta dal sole \\
Ombra non è consistenza, è vanità invece \t è il riflesso (anche montale con lo scalcinato muro) \t trend che mostra assenza della forza ontologica della persona, che fa parte delle macerie \t rimane solo il riflesso della persona, nell acqua e con l'ombra \\
L'acqua ha movimenti impercettibili \t tremula \\
Ombra franta viene cullata nel suo essere inafferrabile, e nell'immagine riflessa della nostra persona che è franta \\
Non c'è speranza nel testo, soltanto nell'immagine dell'immensita \\
Il termine franta descrive l'immagine di questi versi \t una sintassi frante e nucleare, per accumulo di immagini \t lirica pregna di senso \\
Ultima edizione è quella che ha + successo e porta le sue scritture filologiche \\
Illuminazione non viene da un dio, ma dalla consapevolezza di questa vanita

\v

In "Ragioni di una poesia" in "Vita di un uomo" parla di questo \t la poesia ci fa istituire un compromesso, tra il nulla della realta e il tutto del nostro animo, che permette di cogliere il dolore del nulla in cui siamo immersi \\
In "vita di un uomo" guida anche alla comprensione dei suoi testi \\
\ss[Ungaretti per i posteri]
Nel contesto della poesia del tempo, ci fa misurare con uno spartiacque \t fino al 30 il ruolo della poesia ha significativtà che non puo prescindere da ungaretti \\
Anche dopo il fenomeno della WW2 \t anche se qua si ha una esigenza + prosastica, che è anche abbinata alla cinematografia

\ss[Sfaciume - Clemente Rebora]
Clemente Rebora \t focalizza il clima devastato di questi anni \\
Questi poeti appartengono alla corrente dei poeti espressionisti vociani \t La Voce era una rivista \t le riviste servono per pubblicizzare le opere, e raccolgono testi e li pubblicano

\v

Qua la parola diventa il simbolo di un ??? \\
Sfaciume è la parola che diventa spazzatura \t come nel "Inno a satana " di Card. c'è parola portata all'estremo \\
Questa evocazione della morale si perde nelle immagini \t abbiamo unva vita che descrive il mondo attuale, con termini che indicano un abbassamento tonale al verso 11, 12 \t polvere e peste v 13 \\
Strofa che va dal 15 al 20, devastazione del paesaggio ungarettiano in contesto ludico (palloncini, giochi dei ragazzi, con teste vuote (e corpo smilzo)) \\
Dal v 25: volgarita dei versi, che descrive la perdita d'aureola \\
Freschezza irrequieta \t di solito è vitalità, qua no \t il poeta e il suo nume si immolano ad essere solo qualcosa di muto \\
La crisi dell'intellettuale ha portato a una voce che crede nell illuminazione di dio (come nel caso di Rebora) \\
La volgarita di questi versi supera quella di altri \t ma come mai? \\
Si trova come quel clima di distruzione totale \\
Quello che accade con i futuristi, qua si vede coma abbassamento della morale \t parole che entrano nella poetica \t Gioanola e Ermengaldo (sullo sfacelo linguistico, essendo filologo) commentano \\
Questo evoca il deterioramento totale che nella dissolutezza dei costumi porta alla ??

\s[Montale]
Si correla con Dante 

\ss[Vita]
Nella sua parabola di vita, nasce nel 1896 e muore nel 1981 (a genova e muore a milano) \t non ha formazione classica \\
Frequenta istituto tecnico, partecipa alla prima guerra mondiale come fante \\
Vive il post dannunzio con dolorosa tolleranza, che lo porta ad associarsi all espressione che per superare d annunzio bisogna attraversarlo \\
Vive infanzia tra genova e monterosso \t paesaggio ligure, con luoghi impervi in cui si trova natura arida, sofferente ed evocativa \\
La sua vocazione non è quella della guerra \t partecipa ma si associa a un gruppo culturale che si esprime con il nome di Govetti \t e non cela la sua passione per la letteratura \\
Letterato per passione, mentre ragioniere per necessita \t la cultura dell istituto tecnico è orientata a una professione \\
Firma il manifesto degli intellettuali antifascti di Croce \t pubblica "Il venticinque" su Svevo \\
"Ossi di seppia" \t raccolta che è il montale \t tutto parte da qua \t è una raccolta che gli da visibilita e successo \\
La sua parabola di vita lo porta a firenze, partecipa al gabinetto Vessiue, conosce gli editori + prestigiosi ma frequenta il bar dove gli intellettuali si introvavano, tra cui Mario Soldati che diventa amico con lui \\
Dovra abbandonare l incarico al gabinetto, perche non si iscrive al partito fascista \t tra il 32 e il 34 pubblica anche "La casa dei doganieri" \\
La seconda opera ha titolo "Le occasioni" \t negli "ossi di seppia" non ci sono possibilità \t sono fragili e facili a rompersi \\
"Le occasioni" \t titolo perche? \t a volte accade
\end{document}
